\section{Identification and Estimation}\label{sec:identification}

Three objects require estimation: the type-specific cost distributions $F_c^k$, the auction-level scale $a_t$, and the equilibrium entry counts. The third is directly observed. The first two require structure---$F_c^k$ from the design of the Preg\~ao auction \citep{haile2003}, $a_t$ from the variance decomposition in Assumption~\ref{a:uh} \citep{krasnokutskaya2011}. Two diagnostic checks discipline the identification: a cross-modality consistency test that compares the Convite GPV recovery against the Preg\~ao drop-out recovery in a stratum unaffected by the policy, and a Turnbull NPMLE that brackets the all-bidders estimator from below by removing the winner-censoring bias. Both are designed to fail visibly if the identifying assumptions are violated; both are passed.

\subsection{Type-specific cost distributions}

Under Assumption~\ref{a:ipv} and the English-reverse interpretation, the drop-out price of each loser is a point observation of that loser's cost. The main-text estimator is an \emph{all-bidders} ECDF that adds the winner's final clock price to the stratum-specific ECDF; in a descending-clock auction this final price is approximately $c_{(2)}$, a tighter upper bound on the winner's cost than the zero imputed by losers-only omission. Two alternatives bracket all-bidders: losers-only overstates $F_c^k$ in the left tail by omitting the lowest cost draw per auction; the Turnbull NPMLE (Section~\ref{sec:robustness}) treats the winner as left-censored at $c_{(2)}$ and delivers point identification, with the strict ordering $F_c^{\text{Turnbull}} \le F_c^{\text{all-bidders}} \le F_c^{\text{losers-only}}$ in the left tail. Across the eight (class $\times$ type $\times$ period) strata (Table~\ref{tab:v3_turnbull_fc}), the ordering holds at the median in all 8 strata and at the 75th percentile in 6 of 8 (the two pharmaceutical non-SME cells reverse by \turnGapLo--\turnGapHi{} of the reference price, within the bootstrap CI). The BNE decomposition computed on the Turnbull input differs from the all-bidders baseline by at most \turnbullDeviationPct~percent of the total effect (Table~\ref{tab:v3_sensitivity_fc}); losers-only is the input the DiD benchmark of \ref{app:did} effectively absorbs through the item fixed effects.

\subsection{Cross-modality consistency}

A single testable hypothesis disciplines the descending-clock specification: under correct asymmetric-IPV, the \citet{gpv2000} inversion of Convite first-price bids and the \citet{haile2003} drop-out point-ID of Preg\~ao exits should recover the same $F_c^k$ in any stratum unaffected by the policy reform. Pharma non-SME Pre is the natural test cell---it predates the cutoff and is staffed by bidders the set-aside does not target. The two recoveries converge to within one percentage point at the median in that stratum (Figure~\ref{fig:v3_cross_modality_uh}, upper-right panel), and remain within the bootstrap CI throughout the support. Two independent structural recoveries of the same primitive land within sampling error of each other---a non-trivial passage of a falsification test, since systematic divergence between FPSB and English-reverse recoveries would have flagged either an IPV violation or the kind of bidder-group coordination \citet{conley2016} tests detect. Residual gaps in the other three strata motivate the auction-level heterogeneity correction below. Asymmetry across types, in turn, is consistent with the data rather than imposed: recovered cell-level medians and 75th-percentile estimates differ by \medCostGapLo--\medCostGapHi{} of the reference price across the (pharma $\times$ type) cells (Table~\ref{tab:v3_pregao_fc_summary}), well outside Monte Carlo sampling noise.

\begin{figure}[!htb]
\centering
\includegraphics[width=0.98\textwidth]{../output/figures/fig_v3_cross_modality_uh.pdf}
\caption[Cross-modality identification check after UH correction.]{\textit{Cross-modality identification check, UH-clean $F_{c_\epsilon}$.} Each panel plots the type-specific cost distribution estimated two ways on the same sample: the \citet{gpv2000} inversion on Convite first-price sealed-bid auctions (solid) and the \citet{haile2003} drop-out point-ID on Preg\~ao descending-clock auctions (dashed). Pre and Post periods are colour-coded. Convergence in the pharma non-SME Pre panel (upper right) supports the IPV restriction in that stratum: two independent structural recoveries on the same primitive land within sampling error of each other. Residual gaps in other panels motivate the auction-level heterogeneity correction in Section~\ref{sec:identification}.}
\label{fig:v3_cross_modality_uh}
\end{figure}

\subsection{Auction-level heterogeneity}

The median-cost gap across modalities in non-pharma---up to \crossModalityGapPp{} percentage points of the reference price---is too large to attribute to sampling noise. Reference-price normalization absorbs the systematic scale variation across items; the residual within-class auction-level variance is what Assumption~\ref{a:uh}'s $a_t$ structures. Two items in the same CADMAT class can differ in ANVISA registration burden, in contract complexity, or in volume-specific cost elasticities in ways the reference price does not fully reflect.

I estimate the variance decomposition by method of moments. Within each stratum $(k, \text{pharma}, \text{period})$, $\sigma_a^{2,k}$ is the between-auction variance of auction-mean log-bids and $\sigma_e^{2,k}$ the residual within-auction variance, with intraclass correlation $\rho^k = \sigma_a^{2,k} / (\sigma_a^{2,k} + \sigma_e^{2,k})$. The Preg\~ao (pharma $\times$ type) cells span \iccPharmaSmeAvg{}--\iccPharmaNonSme{} (Table~\ref{tab:v3_uh_variance}); restricted to the three cells outside the low-$n$ pharma-SME-Pre stratum, the range \iccNpSme--\iccPharmaNonSme{} brackets \citet{krasnokutskaya2011}'s highway-procurement benchmark of \krasICC. A best-linear-predictor shrinkage estimates each auction's specific heterogeneity as $\hat{a}_t = \hat\sigma_a^2 N_t / (\hat\sigma_a^2 N_t + \hat\sigma_e^2) \cdot (\bar y_t - \bar y)$, and the UH-clean residuals $\hat e_{it} = y_{it} - \hat a_t$ deliver cleaned cost estimates $c_\epsilon^{\text{clean}} = \exp(\hat e_{it})$ that enter all downstream structural estimation.

The correction closes \uhCloseGapPharmaPct~percent of the cross-modality median gap in pharma SME Post and \uhCloseGapNpPct~percent in pharma non-SME Post; the residual non-pharma gap (approximately \uhResidGapNp{} percentage points after UH cleaning) is what the sample filter $c_\epsilon \le \filterCepsUpper$ partially absorbs. Auction-level heterogeneity is therefore the binding constraint in pharma but only part of the story in non-pharma, where richer item-level heterogeneity presumably operates. The Krasnokutskaya shrinkage assumes Gaussianity of $a_t$ and $e_{it}$; the Turnbull NPMLE in Section~\ref{sec:robustness} dispenses with this assumption and ratifies the estimates.

\subsection{Entry and counterfactual simulation}

Entry is modeled as Poisson arrivals with stratum-specific rates. Observed Pre-period average counts of SME and non-SME bidders per auction are the status-quo arrival rates; Post-period SME counts are equilibrium arrival under the restricted regime; no structural entry cost is jointly estimated with the cost primitives (the reduced-form treatment in the \citet{atheyseira2011} spirit). In the main specification, $F_c^{\text{SME,Post}}$ is treated as an equilibrium-selection object generated by the protected regime rather than forced to equal $F_c^{\text{SME,Pre}}$; Section~\ref{sec:robustness} reports the stricter limiting case. A zero-profit calculation delivers implied entry costs in the \entryCostMin--\entryCostMax{} range (Table~\ref{tab:v3_entry_cost}), with non-SMEs facing roughly \entryCostRatioNp{} times the SME entry cost---consistent with the documentation and certification burden of responding to a BEC tender, heavier for larger firms with broader product portfolios subject to ANVISA regulatory scrutiny. With $F_c^k$ and entry counts in hand, the BNE counterfactual is a Monte Carlo over $\monteCarloDraws$ simulated auctions per (class $\times$ scenario): draw type-specific Poisson counts, sample costs from $F_c^k$, record $c_{(1)}$ (production cost) and $c_{(2)}$ (expected transaction price under revenue equivalence). Standard errors come from a cluster bootstrap at the auction level with $B_{bs} = 500$ replicates.

\subsection{Falsifiability}\label{sec:falsifiability}

Three independently verifiable diagnostics discipline the identification, each capable of falsifying rather than merely qualifying it. The cross-modality recovery in pharma non-SME Pre converges within one percentage point at the median (Figure~\ref{fig:v3_cross_modality_uh}); divergence beyond the bootstrap CI would have flagged either an IPV violation or systematic FPSB-vs-English-reverse mismatch. The estimated ICCs $\rho^k$ across the three Preg\~ao non-anomalous cells (\iccNpSme--\iccPharmaNonSme) bracket \citet{krasnokutskaya2011}'s highway-procurement benchmark of \krasICC; an extreme reading would cast doubt on the multiplicative UH form of Assumption~\ref{a:uh}. The Turnbull NPMLE ordering Turnbull $\le$ all-bidders $\le$ losers-only holds at the median in all 8 strata and at the 75th percentile in 6 of 8 (Table~\ref{tab:v3_turnbull_fc}); a reversal beyond Monte Carlo noise would indicate that winner-censoring is not the only bias separating the three estimators.
